\documentclass[aps,prd,reprint,superscriptaddress,showpacs,floatfix,longbibliography]{revtex4-2}

\usepackage[utf8]{inputenc}
\usepackage[T1]{fontenc}
\usepackage{graphicx}
\usepackage{amsmath}
\usepackage{amssymb}
\usepackage{bm}
\usepackage{hyperref}
\usepackage{xcolor}
\usepackage{booktabs}
\usepackage[protrusion=true,expansion=false]{microtype}
\usepackage[switch]{lineno}

\graphicspath{{figures/}}
\setlength\linenumbersep{4pt}
\raggedbottom

\begin{document}
\linenumbers

\title{$\Sigma$-Gravity: A Coherence-Motivated Empirical Response Tested in Galaxies and Clusters}

\author{Leonard Speiser}
\email[Contact author: ]{leonard@horizon3.net}
\homepage[ORCID: ]{https://orcid.org/0009-0008-8797-2457}
\affiliation{Horizon 3, Independent Research, Los Altos, California, USA}

\date{July 25, 2026}

\begin{abstract}
The dynamics of galaxies and galaxy clusters exceed the predictions obtained from directly observed baryonic matter, a discrepancy conventionally modeled with nonbaryonic dark matter or, at galaxy scales, phenomenological modifications such as modified Newtonian dynamics (MOND). We investigate $\Sigma$-Gravity as an additional nonrelativistic phenomenological approach to selected missing-mass observations, not as a disproof of dark matter or a cosmological replacement for $\Lambda$CDM. Its empirical response is $\Sigma(g_N,B)=1+B h(g_N)$, where $h(g_N)=\sqrt{g^\dagger/g_N}\,g^\dagger/(g^\dagger+g_N)$ and $g^\dagger=cH_0/(4\sqrt{\pi})$; the locked galaxy implementation supplements this acceleration dependence with a bounded factor motivated by rotational support. For a disk-dominated sample of 164 SPARC galaxies, using fixed stellar mass-to-light ratios and no per-galaxy halo fitting, the mean velocity root-mean-square residual is $16.366\ {\rm km\,s^{-1}}$ for $\Sigma$-Gravity and $16.056\ {\rm km\,s^{-1}}$ for the MOND prescription tested with the same baryonic inputs. The paired analysis does not resolve a statistically significant difference between these prescriptions under the tested assumptions. Relative to the acceleration-only form, the bounded factor improves the mean RMS by $0.517\ {\rm km\,s^{-1}}$, with a galaxy-bootstrap 95\% interval of $[0.242,0.795]\ {\rm km\,s^{-1}}$. Thus the factor makes a measurable contribution within the locked predictor, although its endogenous construction does not establish independently measured coherence. The cluster exercise is secondary and illustrative: 42 Fox et al. strong-lensing systems determine one effective amplitude under a simplified baryon proxy, while a no-refit check using 73 radial measurements in 17 non-overlapping CLASH clusters gives a median predicted-to-observed ratio of 1.32 and reveals a radius-dependent bias. The Fox normalization therefore does not transfer successfully as a universal cluster law under the present assumptions. An auxiliary QUMOND action reproduces the response for independently prescribed, spatially constant $B$, but not the endogenous galaxy prescription. Together, these results identify a fixed-form galaxy response whose aggregate difference from the tested MOND prescription is unresolved and expose a falsifiable limitation of its fixed cluster amplitude. Kinematic coherence remains the principal motivating hypothesis for the physical origin of $B$, with independently measured phase-space order providing a defined next test; it is not established by the present analysis. No relativistic or cosmological extension is developed here.
\end{abstract}

\maketitle

\noindent\textit{Manuscript ID 1866133; 4,302 words; 4 figures; 2 tables.}

\section{Introduction}

\subsection{The Missing Mass Problem}

Galaxy rotation curves, galaxy-cluster dynamics, gravitational lensing, the cosmic microwave background, and large-scale structure collectively motivate the standard $\Lambda$CDM framework \cite{ref1,ref2}. Within that framework, nonbaryonic cold dark matter supplies the dominant gravitating mass in galaxies and clusters and plays a central role in cosmological structure formation. The microscopic identity of dark matter nevertheless remains unknown, and individual galaxy analyses commonly require system-specific halo parameters. These facts do not negate the cosmological evidence for dark matter, but they motivate tests of compact empirical relationships between baryonic structure and the observed gravitational discrepancy.

\subsection{Modified Gravity Approaches}

MOND provides the best-known alternative phenomenology at galaxy scales \cite{ref3,ref4}. It relates effective and Newtonian accelerations through a universal acceleration scale and organizes many disk-galaxy observations, including the radial acceleration relation \cite{ref5}. Residual mass discrepancies in galaxy clusters and the need for additional structure in relativistic and cosmological realizations motivate the investigation of other phenomenological responses.

\subsection{$\Sigma$-Gravity: Empirical Response and Coherence Hypothesis}

$\Sigma$-Gravity is evaluated here as an additional phenomenological description of selected missing-mass observations alongside dark-matter and MOND descriptions. Its galaxy implementation combines a fixed acceleration response with a bounded factor motivated by rotational support and uses no per-galaxy response amplitude. For a disk rotation curve, this prescription can be compared directly with a MOND interpolation function or a per-galaxy dark-matter halo fit. The term ``alternative'' refers only to this empirical galaxy-scale choice; it does not imply that dark matter has been disproved or that its cosmological role has been replaced. Relativistic and cosmological extensions lie outside the present scope.

The observations analyzed here identify the combined response amplitude $B$, rather than independently separating an amplitude and a physical coherence variable. The implemented galaxy factor is calculated from the model-predicted velocity rather than independent phase-space measurements, and assigning every cluster the same nominal path length identifies one cluster amplitude rather than a length exponent. Coherence remains the principal physical hypothesis motivating the bounded factor, but it is not an established ingredient of the empirical law. The present calculation is an initial test of whether such a factor contributes within the predictor; future work must determine whether independently measured phase-space coherence, source geometry, matter currents, or another property of baryonic organization predicts $B$ on held-out systems.

\subsection{Relation to Previous Work}

This study defines the empirical response and its identifiable parameters, measures the contribution of the bounded galaxy factor through an acceleration-only ablation, gives an auxiliary fixed-$B$ QUMOND construction, compares the locked galaxy predictor with MOND under common baryonic assumptions, and distinguishes cluster calibration from a no-refit external transfer-failure test. The cluster analysis is retained to expose a limitation, not as a successful application. These analyses preserve a falsifiable phenomenology while specifying what a deeper theory would need to derive.

\subsection{Paper Organization}

Section II defines the empirical response, its identifiable parameters, and an auxiliary fixed-$B$ action. Section III describes the SPARC, cluster, and counterrotation analyses. Section IV reports the results while separating calibration from no-refit evaluation. Section V discusses interpretation and scope, and Section VI summarizes the conclusions.

\medskip\hrule\medskip

\section{Theoretical Framework}

\subsection{QUMOND-Like Field Equations}

Let $\Phi_N$ be the Newtonian potential of the baryonic density $\rho_b$,
\begin{equation}
\nabla^2\Phi_N=4\pi G\rho_b,
\label{eq:poissonN}
\end{equation}
and define $g_N=|\nabla\Phi_N|$. The empirical enhancement is
\begin{equation}
\Sigma(g_N,B)=1+B\,h(g_N),
\label{eq:sigma}
\end{equation}
where
\begin{equation}
\begin{aligned}
h(g_N)&=\sqrt{\frac{g^\dagger}{g_N}}
\frac{g^\dagger}{g^\dagger+g_N},\\
g^\dagger&=\frac{cH_0}{4\sqrt{\pi}}
\simeq9.60\times10^{-11}\ {\rm m\,s^{-2}} .
\end{aligned}
\label{eq:h}
\end{equation}
For $g_N\ll g^\dagger$, $h\sim\sqrt{g^\dagger/g_N}$; for $g_N\gg g^\dagger$, $h\rightarrow0$. The exact transition function is a design choice rather than a first-principles result: the square-root factor supplies the deep-acceleration scaling associated with flat rotation curves and the baryonic Tully--Fisher relation, while the rational gate suppresses the response at high acceleration. The scale $g^\dagger\sim cH_0$ places the transition near the empirical galaxy acceleration scale; the geometric factor $4\sqrt{\pi}$ is retained as a fixed phenomenological choice, not an independently derived constant. This study tests the function and its limiting behavior but does not claim that it is unique among nearby transition laws.

\subsection{Identifiability}

In the deep-acceleration point-mass limit,
\begin{equation}
V^4\rightarrow B^2GM_b g^\dagger .
\label{eq:btfr}
\end{equation}
The baryonic Tully--Fisher normalization therefore constrains $B^2g^\dagger$, not $B$ and $g^\dagger$ independently. The observations considered here identify only the combined amplitude $B$; separate interpretations in terms of $B=A\mathcal C$ are therefore not assumed, and possible meanings of $A$ and $\mathcal C$ are deferred to the Discussion.

\subsection{Galaxy Implementation and Status of Coherence}

The locked galaxy implementation uses the bounded fixed-point factor
\begin{equation}
F(V)=\frac{V^2}{V^2+\sigma^2},
\qquad \sigma=20\ {\rm km\,s^{-1}},
\label{eq:F}
\end{equation}
with
\begin{equation}
\begin{aligned}
B_{\rm gal}&=A_0F(V_\Sigma),&
A_0&=e^{1/(2\pi)},\\
V_\Sigma&=V_{\rm bar}\sqrt{1+B_{\rm gal}h(g_N)}.&&
\end{aligned}
\label{eq:galaxy}
\end{equation}
The form $F=V^2/(V^2+\sigma^2)$ is the simplest bounded ratio suggested by an ordered velocity relative to a second velocity moment. The normalization $A_0=e^{1/(2\pi)}$ is motivated by a mode-counting heuristic, and $\sigma=20\ {\rm km\,s^{-1}}$ is adopted as a characteristic cold-disk dispersion scale. These motivations do not constitute first-principles derivations, and the parameters are treated here as fixed model choices.

The formula was developed with SPARC as its principal galaxy benchmark. Here ``locked'' means that Eqs.~(\ref{eq:sigma}), (\ref{eq:h}), (\ref{eq:F}), and (\ref{eq:galaxy}) and their constants define the central configuration, that this configuration was not selected from the sensitivity grid, and that no per-galaxy response parameter is fitted. SPARC is therefore a benchmark evaluation sample, not a prospectively registered or external validation sample. The scale-length, nuisance, ablation, and weighting calculations are sensitivity tests of the same fixed predictor.

The iteration uses $V_\Sigma$, not the observed velocity, and therefore does not directly fit the observed outcome. However, $F(V_\Sigma)$ is endogenous to the prediction rather than an independently measured phase-space quantity. It is treated as a phenomenological regularizer, not as a measured covariant coherence scalar. Writing $x=V_\Sigma^2$, $b=V_{\rm bar}^2$, and $K=A_0h(g_N)$ reduces the fixed point to $x^2+[\sigma^2-b(1+K)]x-b\sigma^2=0$. The product of its roots is negative, so there is exactly one positive physical solution. Initialization at $V_{\rm bar}$ reached the $10^{-6}\ {\rm km\,s^{-1}}$ tolerance for all 2,745 retained points within 26 iterations; no alternative branch is selected by initialization.

The locked predictor in Eq.~(\ref{eq:galaxy}) does not use the catalog photometric scale length $R_d$. As a separate no-refit structural sensitivity, we test
\[
W(r,R_d)=\frac{r}{R_d/(2\pi)+r},\qquad B_{R_d}=A_0W(r,R_d),
\]
using each galaxy's catalog $R_d$ while holding all other choices fixed. This is the previously specified rational window associated with the model's $R_d/(2\pi)$ construction; no alternate windows were screened in this sensitivity calculation. This candidate replaces $F(V_\Sigma)$; it is not multiplied by it and does not alter the central locked result. The reference length $L_0$ and exponent $n$ do not enter either galaxy calculation. A second velocity moment can schematically be decomposed in a stationary axisymmetric system as $\langle v^2\rangle\simeq v_{\rm rot}^2+\sigma^2$, but that frame-dependent statement is not a covariant derivation of Eq.~(\ref{eq:F}). No general covariant reduction is claimed.

\subsection{Auxiliary Fixed-$B$ QUMOND Action}

An auxiliary action can be written in the QUMOND framework \cite{ref6} when $B$ is independently prescribed and spatially constant. Define
\begin{equation}
z=\frac{|\nabla\Phi_N|^2}{(g^\dagger)^2}
\label{eq:z}
\end{equation}
and
\begin{equation}
Q(z,B)=z+4B\left[z^{1/4}-\arctan\!\left(z^{1/4}\right)\right].
\label{eq:Q}
\end{equation}
Then
\begin{equation}
Q_z=1+B\frac{z^{-1/4}}{1+\sqrt z}
=1+B h(g_N).
\label{eq:Qz}
\end{equation}
The gravitational action density
\begin{equation}
\mathcal L_g=-\frac{1}{8\pi G}
\left[2\nabla\Phi\cdot\nabla\Phi_N-(g^\dagger)^2Q(z,B)\right],
\label{eq:action}
\end{equation}
together with the matter coupling $-\rho_b\Phi$, gives by variation
\begin{equation}
\nabla^2\Phi_N=4\pi G\rho_b,
\qquad
\nabla^2\Phi=\nabla\cdot\left[Q_z\nabla\Phi_N\right].
\label{eq:fields}
\end{equation}
This auxiliary construction reproduces the nonrelativistic response for prescribed, spatially constant $B$. Equation~(\ref{eq:action}) is not an action for the endogenous rotational-support prescription in Eq.~(\ref{eq:galaxy}), does not derive $B$, and does not establish conservation or dynamical consistency for that prescription. If $B$ is promoted to a dynamical field, its kinetic and source terms, Euler--Lagrange equation, and backreaction are required. Noether conservation follows for a closed action with the relevant symmetries and all dynamical fields varied; it cannot be inferred by inserting a value of $B$ from observed or predicted kinematics after variation.

\subsection{Algebraic Approximation}

The rotation-curve calculation uses
\begin{equation}
g_{\rm eff}\simeq\Sigma g_N,\qquad
V_\Sigma\simeq V_{\rm bar}\sqrt{\Sigma}.
\label{eq:algebraic}
\end{equation}
This relation is not exact for a general nonspherical mass distribution because Eq.~(\ref{eq:fields}) contains a curl-field contribution. We compare the algebraic acceleration with numerical solutions of the fixed-$B$ QUMOND equation for three analytic axisymmetric disk reconstructions. The diagnostic uses spatially constant $B=1$ and does not recompute the endogenous factor from the numerical circular velocity. A three-dimensional periodic FFT solver uses a $65^3$ grid, a box half-width of $8R_d$, and an exponential--${\rm sech}^2$ disk with scale height $0.2R_d$; the comparison spans $0.75\le r/R_d\le5$. A $49^3$-to-$65^3$ check gives 0.72\% median and 3.67\% maximum acceleration change. The comparison quantifies fixed-response geometry error; it is not a validation of the endogenous prescription or exact SPARC geometries.

\subsection{Path Length, Lensing, and Compact Sources}

Writing the Fox response through a path-length power law assigns every system $L=600$ kpc. Consequently that sample identifies one effective cluster amplitude, not the exponent of a general path-length functional. No unique operational definition of $L$ for an arbitrary three-dimensional baryonic distribution is specified here, so $A(L)$ is not used as a predictive cross-system relation; path length is retained only as a hypothesis. The previously associated value $n=0.27$ is only a reparameterization of the Fox-calibrated amplitude for adopted reference constants, not an independently measured law or part of the canonical response.

The nonrelativistic action determines neither photon propagation nor the two relativistic metric potentials. Equality between a phenomenological dynamical mass and a lensing mass is therefore a closure assumption, not a derived prediction. At Solar-System accelerations $h(g_N)$ is small, but this is not a parameterized post-Newtonian calculation and does not derive the assumption that a compact source has no additional response. Equivalence-principle, post-Newtonian, lensing-slip, and gravitational-wave behavior remain open.

\begin{table*}[t]
\caption{\label{tab:roles}Role of each dataset and analysis.}
\begin{ruledtabular}
\begin{tabular}{llll}
Dataset or analysis & Role & Independent unit & Parameter role\\
\hline
SPARC disk sample & Benchmark evaluation & Galaxy & No per-galaxy parameter\\
SPARC scale-length test & Structural sensitivity & Galaxy & No fitted parameter\\
Fox clusters & Illustrative in-sample calibration & Cluster & One cluster amplitude\\
Repeated Fox splits & Calibration-stability diagnostic & Cluster within Fox & Refits each training subset\\
Non-overlapping Tian/CLASH profiles & No-refit profile check & Cluster, radii grouped & No\\
Matched MaNGA/JAM catalog & Secondary diagnostic & Galaxy/control set & No\\
Numerical QUMOND disks & Approximation diagnostic & Reconstructed disk & No\\
\end{tabular}
\end{ruledtabular}
\end{table*}

\begin{table*}[t]
\caption{\label{tab:parameters}Parameter and assumption accounting.}
\begin{ruledtabular}
\begin{tabular}{llll}
Quantity & Role & Status & Principal limitation\\
\hline
$g^\dagger$ & Acceleration scale in $h(g_N)$ & Fixed choice & BTFR constrains $B^2g^\dagger$\\
$A_0=e^{1/(2\pi)}$ & Galaxy normalization before $F$ & Fixed choice & Not uniquely derived from data\\
$\sigma=20\ {\rm km\,s^{-1}}$ & Regulator in $F(V_\Sigma)$ & Fixed & Endogenous; sensitivity tested\\
$\Upsilon_{\rm disk},\Upsilon_{\rm bulge}$ & Stellar contributions & Fixed & Affect relative performance\\
$B_{\rm Fox}=8.446$ & Cluster response & Illustrative calibration & Not a universal radial amplitude\\
$0.4\times0.15M_{500}$ & Fox aperture baryon mass & Approximation & Material sensitivity\\
Lensing closure & Connects response to lensing target & Assumed & Slip is undetermined\\
$R_d$ & Photometric scale length & Secondary no-refit test & Does not beat assignment controls\\
$L_0,n$ & System-scale path hypothesis & Not used in galaxy result & No operational 3D functional\\
\end{tabular}
\end{ruledtabular}
\end{table*}

\medskip\hrule\medskip

\section{Data and Methodology}

\subsection{SPARC Galaxy Sample}

SPARC provides near-infrared photometry and resolved H\,I/H$\alpha$ rotation curves for 175 disk galaxies \cite{ref7}. The local archive contains 171 usable rotation-curve files. Baryonic velocities are constructed with fixed mass-to-light ratios
\begin{equation}
\Upsilon_{\rm disk}=0.5,\qquad
\Upsilon_{\rm bulge}=0.7,
\label{eq:ml}
\end{equation}
and no galaxy-specific halo or response amplitude is fitted. Because the circular-speed approximation is least secure where a spheroidal bulge dominates, a radial point is excluded from the primary disk analysis when the scaled bulge contribution is at least 30\% of $V_{\rm bar}^2$. Galaxies with fewer than three remaining points are excluded. The locked sample contains 164 galaxies and 2,745 disk-dominated radial points. Sensitivity checks use 20\% and 40\% thresholds and all 171 usable galaxies with 3,373 valid points.

\subsection{MOND Comparison and Statistics}

MOND is evaluated with the same baryonic inputs using
\begin{equation}
\begin{aligned}
\nu(x)&=\frac{1}{1-e^{-\sqrt{x}}},&
x&=\frac{g_N}{a_0},\\
a_0&=1.2\times10^{-10}\ {\rm m\,s^{-2}}.&&
\end{aligned}
\label{eq:mond}
\end{equation}
For each galaxy, we calculate an unweighted velocity RMS residual for $\Sigma$-Gravity and MOND. This gives each retained radius equal weight within a galaxy and each galaxy equal weight in the sample. In the secondary uncertainty-weighted calculation, $w_{ij}=1/\sigma_{V,ij}^2$ and the within-galaxy statistic is $[\sum_jw_{ij}(V_{{\rm model},ij}-V_{{\rm obs},ij})^2/\sum_jw_{ij}]^{1/2}$; these galaxy-level values are then averaged with equal weight across galaxies. Galaxies, not radial measurements, are the independent resampling units. Uncertainty intervals use 20,000 galaxy-bootstrap resamples. We use an exact binomial comparison of the fraction of galaxies for which each model has lower RMS, with no ties in the primary result and null probability 0.5, and a paired sign-flip test of the mean RMS difference. The scale-length sensitivity uses catalog $R_d$ without fitting, plus 2,000 galaxy-level assignment permutations and a fixed-median negative control.

The frozen nuisance grid contains 81 combinations of common disk and bulge mass-to-light ratios, a global distance scale, a global inclination offset, and the $A_0$ normalization. A separate diagnostic varies $\sigma$ from 10 to $50\ {\rm km\,s^{-1}}$. These are sensitivity analyses, not fitted alternatives.

\subsection{Galaxy Cluster Samples}

The Fox et al. catalog relates strong-lensing strength to cluster properties \cite{ref8}. The sample used here contains 42 clusters with spectroscopic redshifts and $M_{500}>2\times10^{14}M_\odot$. The baryon approximation is
\begin{equation}
M_b(<200\ {\rm kpc})=0.4\times0.15\,M_{500}.
\label{eq:clusterbaryon}
\end{equation}
The factor 0.4 is a simplified baryon-concentration prescription rather than a measured gas-plus-stellar profile and provides a reproducible baseline for the Fox calibration. The calibrated cluster amplitude is $B_{\rm Fox}=8.446$. Writing this amplitude through a power law with $L=600$ kpc for every cluster does not identify a path-length exponent because $L$ does not vary within the sample. In particular, the previously associated value $n=0.27$ is only a reparameterization of this single calibrated amplitude for adopted reference constants; it is not used in the canonical response and is not independently measured by the cluster data. Repeated train/test splits within Fox refit the same catalog and are described only as calibration-stability checks.

The no-refit radial check uses the Tian et al. CLASH radial-acceleration catalog \cite{ref9}. Cluster identifiers were case-normalized, stripped of non-alphanumeric characters, standardized for the documented Abell and MACS aliases, and compared for exact equality with a Fox calibration name. This procedure excluded MACS0416, MACS0717, and MACS1149, leaving 73 measurements in 17 clusters. The Fox-calibrated $B_{\rm Fox}=8.446$, $h(g_N)$, and $g^\dagger$ are fixed before comparison. For the trend in $r_i=\log_{10}g_{{\rm pred},i}-\log_{10}g_{{\rm tot},i}$, the propagated uncertainty is $\sigma_{r,i}^2=\sigma_{\log g_{\rm tot},i}^2+[s_i\sigma_{\log g_{\rm bar},i}]^2$, where $s_i=d\log_{10}g_{\rm pred}/d\log_{10}g_{\rm bar}$, and the regression uses weights $1/\sigma_{r,i}^2$. Repeated radii are retained within cluster, and complete clusters---not individual radii---are the resampling units in 5,000 bootstrap draws.

\subsection{Counterrotation Diagnostic}

The counterrotator catalog is based on the MaNGA sample of Bevacqua et al. \cite{ref10}, and the secondary dynamical quantities are drawn from MaNGA DynPop \cite{ref11}. Repeat observations are reduced to one row per galaxy by retaining the lowest-JAM-$\chi^2$ entry. Sixty-two counterrotating systems with complete fields are matched to five unique controls each on stellar mass, physical size, S\'ersic index, axis ratio, inclination, redshift, and JAM fit quality. The greedy Euclidean nearest-neighbor procedure standardizes covariates on the combined pool, processes the hardest case first, uses no replacement and no caliper, and retains all 62 eligible cases. Balance is assessed with the absolute standardized mean difference (SMD), using 0.1 as the stated balance threshold. Bootstrap samples resample the 62 complete matched sets while retaining each five-control mean.

The compared outcome, $f_{\rm DM}(<R_e)$, is inferred within a Newtonian JAM/NFW model. It is not a direct observable and is retained only as a secondary diagnostic. A direct test would require common forward modeling of MaNGA velocity and dispersion maps with complete galaxies held out.

\medskip\hrule\medskip

\section{Results}

\subsection{SPARC Galaxy Rotation Curves}

For 164 galaxies and 2,745 disk-dominated radial points, the mean galaxy RMS is $16.366\ {\rm km\,s^{-1}}$ for $\Sigma$-Gravity and $16.056\ {\rm km\,s^{-1}}$ for MOND. Define the paired galaxy contrast
\begin{equation}
\Delta_i={\rm RMS}_{\Sigma,i}-{\rm RMS}_{{\rm MOND},i}.
\label{eq:delta}
\end{equation}
Its galaxy-weighted mean is $+0.309\ {\rm km\,s^{-1}}$, with bootstrap 95\% interval $[-0.040,+0.659]\ {\rm km\,s^{-1}}$. $\Sigma$-Gravity has lower RMS for 71 of 164 galaxies, or 43.3\%. There are no exact ties; the two-sided exact binomial test uses null probability 0.5 and gives $p=0.101$. The paired sign-flip value is $p=0.088$. The analysis therefore does not resolve a statistically significant difference between the prescriptions under the tested assumptions; no equivalence or noninferiority claim is made.

Across the 81 frozen nuisance combinations, including common changes in $A_0$, the mean paired contrast ranges from $-3.286$ to $+2.779\ {\rm km\,s^{-1}}$ and retains the central sign in 64.2\% of cases. Changing $\sigma$ from the central 20 km s$^{-1}$ to 10, 30, and 50 km s$^{-1}$ gives mean $\Sigma$-Gravity RMS values of 16.695, 16.347, and 17.600 km s$^{-1}$, respectively, compared with 16.366 km s$^{-1}$ at 20 km s$^{-1}$. The 30 km s$^{-1}$ value is only 0.019 km s$^{-1}$ lower than the central choice, and no alternative is selected. All four $\sigma$ choices retain a positive mean $\Sigma$-minus-MOND contrast. These sensitivities show that the fixed choices are not uniquely identified. Replacing Eq.~(\ref{eq:galaxy}) with the acceleration-only form $B=A_0$ gives a mean RMS of $16.882\ {\rm km\,s^{-1}}$. The locked endogenous factor improves the mean RMS by $0.517\ {\rm km\,s^{-1}}$, with bootstrap 95\% interval $[0.242,0.795]\ {\rm km\,s^{-1}}$. This is an internal ablation of the empirical equation, not evidence for an independently measured coherence effect.

Uncertainty weighting gives mean RMS values of $15.847$ and $15.654\ {\rm km\,s^{-1}}$ for $\Sigma$-Gravity and MOND, respectively, with paired mean $+0.193\ {\rm km\,s^{-1}}$ and 95\% interval $[-0.177,+0.564]\ {\rm km\,s^{-1}}$. The mean $\Sigma$-minus-MOND contrast remains positive at 20\%, 30\%, and 40\% bulge thresholds and in the all-valid-points sample. These checks likewise do not resolve a statistically significant aggregate difference.

\begin{figure*}[t]
\centering
\includegraphics[width=\textwidth]{figure_1_sparc_paired.pdf}
\caption{\label{fig:sparc}Locked SPARC comparison and nuisance sensitivity. Left: per-galaxy velocity RMS for $\Sigma$-Gravity and the tested MOND prescription. Center: distribution of the paired contrast; its 95\% galaxy-bootstrap interval includes zero. Right: mean contrast for the 81 frozen nuisance combinations. Negative values favor $\Sigma$-Gravity and positive values favor MOND. Neither the primary paired interval nor this sensitivity grid establishes a statistically significant difference or equivalence.}
\end{figure*}

\subsection{Photometric Scale-Length Hypothesis}

The catalog-scale-length candidate gives mean galaxy RMS $16.513\ {\rm km\,s^{-1}}$, compared with $16.366\ {\rm km\,s^{-1}}$ for the locked predictor. Their paired difference is $+0.147\ {\rm km\,s^{-1}}$, with bootstrap 95\% interval $[-0.088,+0.385]\ {\rm km\,s^{-1}}$; this does not resolve a statistically significant paired difference and is not an equivalence test. The candidate improves on the acceleration-only ablation by $0.370\ {\rm km\,s^{-1}}$ but is worse than the tested MOND prescription by $0.456\ {\rm km\,s^{-1}}$.

The actual catalog assignments do not outperform random reassignment (one-sided $p=0.711$), and a common median $R_d=2.3$ kpc gives a lower mean RMS of $16.413\ {\rm km\,s^{-1}}$. Thus radial suppression is useful relative to the acceleration-only response, but this test does not support measured photometric scale length as its controlling variable. Equation~(\ref{eq:galaxy}) is retained as the locked predictor.

\subsection{Illustrative Galaxy Cluster Calibration and No-Refit Check}

Equation~(\ref{eq:clusterbaryon}) is used to calibrate the common cluster response on the 42 Fox systems, giving a median predicted-to-observed mass ratio of 0.987 and a scatter of 0.132 dex. Because the amplitude and simplified baryon proxy are inseparable in this calculation, the exercise is retained only as an illustrative in-sample calibration. Varying the 0.4 concentration factor by 25\% changes the predicted mass ratio by approximately 30\%, so the calibrated amplitude cannot be separated from the assumed baryonic aperture mass.

With $B_{\rm Fox}=8.446$ held fixed, the 73 measurements in 17 non-overlapping CLASH clusters give a median predicted-to-observed ratio of 1.318 and an RMS residual of 0.188 dex. Disjoint-sample median ratios are 1.17 at 100 kpc, 1.23 at 200 kpc, 1.61 at 400 kpc, and 1.96 at 600 kpc. The weighted log-residual slope is 0.162 dex per dex in radius, with a 5,000-draw cluster-bootstrap 95\% interval $[0.115,0.217]$. Thus the Fox-calibrated amplitude increasingly overpredicts the profile at larger radius under the present baryonic and lensing assumptions. This no-refit result does not identify a replacement amplitude or cluster formula; it shows that the Fox-calibrated amplitude does not transfer successfully to the CLASH radial data and provides no support for a universal fixed cluster-amplitude law under the present assumptions.

\begin{figure*}[t]
\centering
\includegraphics[width=\textwidth]{figure_2_cluster_roles.pdf}
\caption{\label{fig:clusters}Illustrative cluster calibration and no-refit radial evaluation. Left: the 42 Fox clusters under the simplified baryon proxy and calibrated $B_{\rm Fox}$; this is illustrative in-sample calibration. Right: 73 Tian/CLASH radial measurements in 17 clusters after the fixed name-normalization rule excludes MACS0416, MACS0717, and MACS1149, with $B_{\rm Fox}$ frozen. Large markers show disjoint-sample radius-bin medians. The systematic increase with radius shows that the fixed Fox-calibrated amplitude does not transfer without bias under the present baryonic and lensing assumptions.}
\end{figure*}

\subsection{Algebraic Approximation}

Numerical three-dimensional QUMOND solutions for analytic axisymmetric disk reconstructions representative of F574-2, UGC05716, and NGC3741 give median absolute acceleration differences of 5.19\%, 4.88\%, and 3.96\%, respectively, relative to Eq.~(\ref{eq:algebraic}). Maximum local differences are 20.54\%, 18.30\%, and 7.73\%. These are fixed-$B$ field-geometry checks, not solutions of the endogenous Eq.~(\ref{eq:galaxy}) model. The radius-dependent geometry error is comparable to or larger than the aggregate RMS separation between $\Sigma$-Gravity and MOND for some galaxies and should be included in future object-level likelihoods.

\begin{figure}[t]
\centering
\includegraphics[width=\columnwidth]{figure_3_qumond_approximation.pdf}
\caption{\label{fig:qumond}Fractional acceleration difference between the algebraic relation and numerical fixed-$B$ QUMOND solutions for three analytic disk reconstructions. The result is an approximation-error diagnostic, not a fit to the locked endogenous prescription or full observed gas and bulge maps.}
\end{figure}

\subsection{Counterrotation}

Matching reduces the maximum absolute SMD from 0.684 to 0.066, below the 0.1 balance threshold. The matched mean difference is
\begin{equation}
\Delta f_{\rm DM}
=f_{{\rm DM},{\rm counter}}-f_{{\rm DM},{\rm control}}
=-0.0081 ,
\label{eq:fdm}
\end{equation}
with matched-set bootstrap 95\% interval $[-0.0577,+0.0453]$. The result is consistent with zero and does not reproduce the large unmatched association. Because $f_{\rm DM}$ is JAM/NFW-derived, this is not a direct test of $\Sigma$-Gravity. Counterrotation remains a proposed future discriminant rather than a confirmed prediction.

\begin{figure}[t]
\centering
\includegraphics[width=\columnwidth]{figure_4_counterrotation_matched.pdf}
\caption{\label{fig:counter}Matched counterrotation diagnostic. Left: covariate balance before and after matching. Right: matched difference in JAM/NFW-derived $f_{\rm DM}(<R_e)$; the matched-set bootstrap 95\% interval includes zero.}
\end{figure}

\medskip\hrule\medskip

\section{Discussion}

\subsection{A Bounded Alternative}

$\Sigma$-Gravity is an additional empirical response law for selected missing-mass observations. At galaxy scale, it can be evaluated alongside a fitted halo or a MOND interpolation function using a specified baryon-dependent prescription. The locked SPARC result shows that this fixed-form prescription remains close to MOND in aggregate without fitting an individual response amplitude to each galaxy. That is an interesting empirical result, but it is not evidence that dark matter is absent or that the model supersedes MOND.

At cluster scale, the implementation is less successful. The Fox exercise is retained only as an illustrative calibration under a simple baryon prescription, not as central positive evidence. The no-refit CLASH profiles show that this normalization does not transfer successfully and provides no support for a universal fixed cluster-amplitude law under the present assumptions. The Fox value is therefore retained only as a calibration baseline; no replacement cluster formula is introduced.

\subsection{Coherence and Path Length as Hypotheses}

The empirical response does not identify its physical mechanism. The principal hypothesis motivating the present construction is that the effective response depends on the organization or dynamical state of the baryonic source. Kinematic coherence, or another correlated property of organized matter, could contribute to a dynamical order parameter. The bootstrap-supported internal improvement of the bounded factor over the acceleration-only ablation supports continued investigation of this possibility within the model, while the endogenous construction prevents it from being interpreted as a detection of physical coherence. For independently measured phase-space data, one operational estimator would be
\begin{equation}
\mathcal C_{\rm kin}=
\frac{|\bar{\bm v}_{\rm stream}|^2}
{|\bar{\bm v}_{\rm stream}|^2+
{\rm tr}\,\bm{\sigma}^2}
\label{eq:ckin}
\end{equation}
in the system barycentric frame. This quantity would be small for a dispersion-supported cluster, whereas the Fox calibration uses a single effective response rather than an independently measured $\mathcal C_{\rm kin}$. One measured coherence definition therefore does not currently explain both disks and clusters.

No unique operational path-length functional for a general three-dimensional baryonic distribution is specified here. Because every Fox cluster receives $L=600$ kpc, the sample cannot identify the exponent $n$ independently; the value $n=0.27$ is not part of the canonical response or an independently measured result. Consequently, $A(L)$ is not used as a predictive cross-system relation. The direct SPARC sensitivity shows that the tested window improves on acceleration-only behavior but that actual catalog $R_d$ assignments carry no advantage over permuted or fixed controls. Coherence, path length, and source organization should advance only if independently measured quantities improve held-out predictions relative to an acceleration-only model.

A related but distinct future direction is Refracted Gravity, in which a density-dependent gravitational permittivity modifies the Poisson equation and can redirect field lines in nonspherical sources \cite{ref12,ref13,ref14}. This literature illustrates how density and source geometry might enter a developed theory, but it does not derive or validate the empirical $\Sigma$ response.

\subsection{Comparison with MOND}

The tested $\Sigma$-Gravity and MOND prescriptions use the same baryonic inputs and no per-galaxy halo fitting. The paired analysis does not resolve a statistically significant aggregate difference under the locked assumptions, and no equivalence or noninferiority claim is made. This statement is narrower than comparing complete theories: neither the nonrelativistic $\Sigma$ response nor the MOND interpolation function used here provides by itself a full relativistic or cosmological model.

\subsection{Lensing, Solar-System, and Cosmological Limitations}

In a general weak-field metric,
\begin{equation}
ds^2=-(1+2\Psi)c^2dt^2+(1-2\Phi)d\bm{x}^2,
\label{eq:metric}
\end{equation}
nonrelativistic dynamics primarily constrains $\Psi$, while lensing depends on $\Phi+\Psi$. The gravitational slip $\eta=\Phi/\Psi$ is not determined by Eqs.~(\ref{eq:action})--(\ref{eq:fields}). The cluster comparisons are therefore empirical mass comparisons conditional on a lensing closure assumption.

At Solar-System accelerations $h(g_N)$ is small. This is an anomalous-acceleration sanity check, not a calculation of the parameterized post-Newtonian parameter $\gamma$, and the assumption that a compact source has no additional response has not been derived.

No cosmological extension is developed here, and the present response law should not be interpreted as an alternative cosmology. Any such extension would need to address the cosmic microwave background spectrum, primordial abundances, perturbation growth, nonlinear structure formation, and cluster-abundance evolution that form central parts of the $\Lambda$CDM evidence base.

\subsection{Staged Future Investigations}

The open limitations define a staged research program in which additional claims depend on independent observables rather than additional formula flexibility. First, MaNGA velocity and dispersion maps should be forward-modeled for counterrotators and balanced controls under common baryonic and nuisance assumptions, with complete galaxies held out. An acceleration-plus-coherence model should advance only if an independently measured phase-space estimator improves galaxy-grouped held-out performance over the acceleration-only response by a prespecified margin and retains a stable sign across morphology and mass.

Second, a disjoint cluster sample should combine radial X-ray gas, brightest-cluster-galaxy, satellite, and intracluster-light profiles with the full lensing covariance fixed before any response amplitude is inferred. A cluster law should be called predictive only if it transfers without refitting across clusters and radii. Candidate density-, concentration-, or geometry-sourced fields should be specified without system-class switches or per-object parameters and compared with a constant-$B$ baseline on complete held-out systems.

Third, the algebraic disk approximation should be replaced in the likelihood by exact axisymmetric or three-dimensional field solutions using observed gas, stellar, and bulge maps. A first-principles nonrelativistic extension would promote $B$ to an independent field with explicit kinetic and source terms, vary all fields in a closed action, and test stability and conservation before interpreting $B$ as coherence or path length. A relativistic completion would then need to derive both metric potentials, photon propagation, post-Newtonian behavior, equivalence-principle limits, and gravitational-wave propagation. Cosmological background and perturbation calculations would be warranted only if these preceding tests succeed. Confirmatory analyses should use frozen samples and covariance models, object-grouped validation, public code and machine-readable outputs, and independent replication.

\medskip\hrule\medskip

\section{Conclusions}

$\Sigma$-Gravity provides a compact acceleration-based phenomenology that can be considered alongside dark-matter and MOND descriptions of selected missing-mass observations. On a locked sample of 164 SPARC galaxies, the paired analysis does not resolve a statistically significant rotation-curve difference from the MOND prescription tested with the same fixed baryonic assumptions and without fitting an individual halo or response amplitude to each galaxy. This is not an equivalence or noninferiority result. Within the $\Sigma$-Gravity predictor, the bounded rotational-support-motivated factor improves the mean RMS over the acceleration-only form by $0.517\ {\rm km\,s^{-1}}$, with a 95\% interval of $[0.242,0.795]\ {\rm km\,s^{-1}}$. This is a statistically supported internal contribution and provides an initial reason to test the coherence hypothesis with independent kinematics, although its endogenous construction does not establish independently measured coherence.

The Fox cluster result is an illustrative in-sample calibration conditional on a simplified baryon model and an assumed relation to lensing mass, not a principal positive test of the framework. A no-refit CLASH profile check reveals increasing radial overprediction. The calibrated $B=8.446$ is therefore not established as a universal cluster value, and a path-length interpretation remains a hypothesis.

An auxiliary nonrelativistic QUMOND action reproduces the response for independently prescribed, spatially constant $B$, but it is not an action for the endogenous galaxy prescription, and the physical origin and dynamics of that amplitude remain unknown. Coherence or another property of baryonic organization is a possible explanation worth testing; it is not established by the current data. The matched counterrotation result is consistent with zero and is reported as a negative secondary diagnostic.

Together, these results document the performance of a fixed-form galaxy response with no per-galaxy amplitude, an auxiliary fixed-$B$ QUMOND construction, and an external cluster test that reveals radius-dependent biased transfer of the Fox-calibrated amplitude under the present assumptions. They motivate further prespecified testing of $\Sigma$-Gravity as a galaxy-scale phenomenological alternative, without claiming to invalidate dark matter or supersede MOND. The next evidentiary step is not a more flexible response law, but a preregistered test using independently measured phase-space order and cluster baryon profiles on held-out systems. Those tests can determine whether coherence or another property of baryonic organization has causal explanatory value, whether the response reflects other new physics, or whether it is a compact description of correlated astrophysical structure; relativistic and cosmological development remains contingent on their outcome.

The measured SPARC scale lengths do not improve the tested radial-window extension over the locked predictor or its assignment controls. This negative result preserves the empirical core while narrowing the claim: coherence, scale length, and source organization are candidate explanations for $B$, not observationally established dependencies.

\medskip\hrule\medskip

\section{Data and Code Availability}

SPARC data are publicly available at \url{http://astroweb.cwru.edu/SPARC/}. The CLASH radial-acceleration catalog is available through VizieR as J/ApJ/896/70. The Fox calibration table, frozen residuals, split definitions, matched samples, parameter diagnostics, and figure-generation code are provided at \url{https://github.com/lrspeiser/sigmagravity/tree/main/Publications/Frontiers} and in the Supplementary Material.

\section{Author Contributions}

LS conceived the study, developed the model, assembled the data and code, performed the analyses, interpreted the results, and wrote and revised the manuscript.

\section{Funding}

The author declares that no financial support was received for the research, authorship, and/or publication of this article.

\section{Conflict of Interest}

The author declares that the research was conducted in the absence of any commercial or financial relationships that could be construed as a potential conflict of interest.

\begin{acknowledgments}
The author thanks Emmanuel N. Saridakis, Rafael Ferraro, and Tiberiu Harko for earlier discussions concerning theoretical consistency and modified-gravity frameworks.
\end{acknowledgments}

\begin{thebibliography}{99}

\bibitem{ref1} F. Zwicky, Helv. Phys. Acta \textbf{6}, 110 (1933).

\bibitem{ref2} Planck Collaboration, Astron. Astrophys. \textbf{641}, A6 (2020), doi:10.1051/0004-6361/201833910.

\bibitem{ref3} M. Milgrom, Astrophys. J. \textbf{270}, 365 (1983), doi:10.1086/161130.

\bibitem{ref4} R. H. Sanders and S. S. McGaugh, Annu. Rev. Astron. Astrophys. \textbf{40}, 263 (2002), doi:10.1146/annurev.astro.40.060401.093923.

\bibitem{ref5} S. S. McGaugh, F. Lelli, and J. M. Schombert, Phys. Rev. Lett. \textbf{117}, 201101 (2016), doi:10.1103/PhysRevLett.117.201101.

\bibitem{ref6} M. Milgrom, Mon. Not. R. Astron. Soc. \textbf{403}, 886 (2010), doi:10.1111/j.1365-2966.2009.16184.x.

\bibitem{ref7} F. Lelli, S. S. McGaugh, and J. M. Schombert, Astron. J. \textbf{152}, 157 (2016), doi:10.3847/0004-6256/152/6/157.

\bibitem{ref8} C. Fox, G. Mahler, K. Sharon, and J. D. Remolina Gonz\'alez, Astrophys. J. \textbf{928}, 87 (2022), doi:10.3847/1538-4357/ac5024.

\bibitem{ref9} Y. Tian, K. Umetsu, C.-M. Ko, M. Donahue, and I.-N. Chiu, Astrophys. J. \textbf{896}, 70 (2020), doi:10.3847/1538-4357/ab8e3d.

\bibitem{ref10} D. Bevacqua, M. Cappellari, and S. Pellegrini, Mon. Not. R. Astron. Soc. \textbf{511}, 139 (2022), doi:10.1093/mnras/stab3732.

\bibitem{ref11} K. Zhu et al., Mon. Not. R. Astron. Soc. \textbf{522}, 6326 (2023), doi:10.1093/mnras/stad1299.

\bibitem{ref12} V. Cesare, A. Diaferio, T. Matsakos, and G. Angus, Astron. Astrophys. \textbf{637}, A70 (2020), doi:10.1051/0004-6361/201935950.

\bibitem{ref13} V. Cesare, A. Diaferio, and T. Matsakos, Astron. Astrophys. \textbf{657}, A133 (2022), doi:10.1051/0004-6361/202140651.

\bibitem{ref14} A. P. Sanna, T. Matsakos, and A. Diaferio, Astron. Astrophys. \textbf{674}, A209 (2023), doi:10.1051/0004-6361/202243553.

\end{thebibliography}

\end{document}
